% !TEX root = ../Main.tex
\chapter{测量误差}

测量误差分两类：
\[
\text{误差}\ \begin{cases}\text{系统误差}\\[2pt]\text{偶然误差}\end{cases}
\]

\section{偶然误差}

偶然误差是\textbf{人为的}，常表现为人的判断力带来的偏差，比如\textbf{读数}时多估或少估一点。

它的关键特征是\textbf{双向性}：有时偏大、有时偏小，没有固定方向。正因如此，办法是\textbf{多次实验取平均值}——一正一负相互抵消，让双向性的效果表现为"$0$"。

\section{系统误差}

系统误差是\textbf{实验原理上就引入}的误差，跟做多少次没关系。例如：

\begin{itemize}
    \item 研究自由下落时，小球必然受到空气阻力；
    \item 用电表测量时，电表本身有内阻。
\end{itemize}

它的特征是\textbf{单向性}：总是朝同一个方向偏，多次取平均也消不掉，只能从原理或仪器上去改进。

\rmkb{
    一句话区分：偶然误差靠"多测几次取平均"压下去（双向性 $\to 0$）；系统误差是原理自带的偏差（单向性），平均没用。
}
